\documentclass[../../main]{subfiles}

\begin{document}
\section{Métodos de Demonstração}
\label{sec:matematica:logica-matematica:metodos-de-demonstracao}

\begin{defn}[Método de Demonstração]
  \label{def:matematica:logica-matematica:metodos-de-demonstracao:metodo-de-demonstracao}

  Um método de demonstração é um procedimento lógico utilizado para
  estabelecer a validade de uma proposição a partir de axiomas,
  definições e proposições previamente demonstradas.
\end{defn}

\begin{defn}[Demonstração Direta]
  \label{def:matematica:logica-matematica:metodos-de-demonstracao:demonstracao-direta}

  Seja uma proposição da forma

  \[
    P \Rightarrow Q.
  \]

  Uma demonstração direta consiste em assumir \(P\) como verdadeira e,
  por meio de uma sequência finita de inferências válidas, concluir \(Q\).
\end{defn}

\begin{ex}
  Se \(n\) é um número par, então \(n^2\) é par.

  Assuma

  \[
    n = 2k,
    \quad
    k \in \mathbb{Z}.
  \]

  Então

  \[
    n^2 = (2k)^2 = 4k^2 = 2(2k^2),
  \]

  logo \(n^2\) é par.
\end{ex}

\begin{defn}[Demonstração por Contraposição]
  \label{def:matematica:logica-matematica:metodos-de-demonstracao:contraposicao}

  Seja uma proposição da forma

  \[
    P \Rightarrow Q.
  \]

  Uma demonstração por contraposição consiste em demonstrar a
  proposição logicamente equivalente

  \[
    \neg Q \Rightarrow \neg P.
  \]
\end{defn}

\begin{prop}[Equivalência da Contraposição]
  \label{prop:matematica:logica-matematica:metodos-de-demonstracao:equivalencia-contraposicao}

  Para quaisquer proposições \(P\) e \(Q\),

  \[
    (P \Rightarrow Q)
    \Leftrightarrow
    (\neg Q \Rightarrow \neg P).
  \]
\end{prop}

\begin{proof}

  Considere a tabela-verdade das duas expressões.

  Verifica-se que possuem os mesmos valores de verdade para todas as
  atribuições possíveis.

  Portanto são logicamente equivalentes.

\end{proof}

\begin{defn}[Demonstração por Contradição]
  \label{def:matematica:logica-matematica:metodos-de-demonstracao:contradicao}

  Seja \(P\) uma proposição.

  Uma demonstração por contradição consiste em assumir
  temporariamente a falsidade de \(P\), isto é,

  \[
    \neg P,
  \]

  e deduzir uma contradição.

  Pelo princípio da não contradição conclui-se então que \(P\) é
  verdadeira.
\end{defn}

\begin{ex}

  A irracionalidade de \(\sqrt{2}\) é tradicionalmente demonstrada por
  contradição.

\end{ex}

\begin{defn}[Demonstração por Casos]
  \label{def:matematica:logica-matematica:metodos-de-demonstracao:casos}

  Uma demonstração por casos consiste em decompor o problema em um
  conjunto finito de situações mutuamente exclusivas e demonstrar a
  proposição em cada uma delas.
\end{defn}

\begin{prop}[Princípio da Demonstração por Casos]
  \label{prop:matematica:logica-matematica:metodos-de-demonstracao:principio-casos}

  Se

  \[
    P_1 \lor P_2 \lor \cdots \lor P_n
  \]

  é verdadeira e

  \[
    P_i \Rightarrow Q
  \]

  para todo \(i\), então \(Q\) é verdadeira.
\end{prop}

\begin{proof}

  Como pelo menos um dos casos \(P_i\) ocorre e cada um deles implica
  \(Q\), segue que \(Q\) é verdadeira.

\end{proof}

\begin{defn}[Indução Matemática]
  \label{def:matematica:logica-matematica:metodos-de-demonstracao:inducao-matematica}

  Seja \(P(n)\) uma proposição definida para todo
  \(n \in \mathbb{N}\).

  O princípio da indução matemática afirma que, se:

  \begin{enumerate}
    \item \(P(0)\) é verdadeira;

    \item para todo \(n \in \mathbb{N}\),

      \[
        P(n)
        \Rightarrow
        P(n+1),
      \]

      então \(P(n)\) é verdadeira para todo
      \(n \in \mathbb{N}\).
  \end{enumerate}

\end{defn}
\end{document}
